\section[PHT3D Exercise 6 Diffusion-controlled transport of radioniclides]
{PHT3D Exercise 6 Diffusion-controlled transport of radionuclides}

One of the most important roles of reactive transport modelling
is its application to assess and predict the transport of radionuclides,
particularly in the context of the safety of the terminal 
storage of nuclear waste materials.  
One of the major stategies for storing
radioactive waste underground is to select geological fromations
that contain a suitable clay layer that will cause the transport
of the radionuclides away from the storage location 
to be strongly retarded. The transport across the clay layer
occurs more or less exclusively through molecular diffusion.
It is therefore assumed to be slow enough to allow for a  
substantial reduction of mass via radioactive decay. 
However the clay layer will typically be in contact 
with the host rock and some mineral transformations
that may impact the transport behavior of the radionuclides 
might occur. 

As an example of such a problem we will in this exercise simulate 
the mineral dissolution/precipitation reactions that occur 
at a kaolinite/granite boundary. Once the model is set up
we will use it to analyse the role of a mineral barrier 
on the diffusion-controlled transport of uranium in such a system.
The general set up of the problem and dimensions 
are schematicall shown in Figure \ref{setup_u_diff}.

Figure 1 : domain extent, porosity and mineral composition 
(Montmor. Stands for Ca-Montmorillonite), some Fe(OH)3(a) 
is present over the whole domain

\begin{figure}
\centering
\includegraphics[width=1.0\textwidth]{figures/u_diff.eps}
\caption[Schematic overview: Model extent, porosity and mineral composition. 
Note, Montmor. refers to Ca-Montmorillonite. Some Fe(OH)3(a) is 
present over the whole domain.]{Schematic overview: Model extent, 
porosity and mineral composition. 
Note, Montmor. refers to Ca-Montmorillonite. Some Fe(OH)3(a) is 
present over the whole domain.}
\label{fig:setup_u_diff}
\end{figure}


\subsection{Model Domain}

\begin{itemize}
\item Build the geometry of the model, as a 2D horizontal model, 
with 6 m length and 1 m width, with dx = 0.05 m and dy = 1m 
(just one cell in y direction)
\item Set flow boundary conditions to -1 on both sides of the model and 
set the initial head values at 1 m over the whole domain.
\item In Parameters/Flow/Parameters, 
in Dis.2 set time to years. 
Your model will be in years
\item In Parameters/Flow/Parameters, 
in Dis.8 set the internal time steps to 10. 
This value also corresponds to the number of chemistry 
calculation time steps that will occur during a flow time step.
\item In Parameters/Model/time, set the final time to 1000 y, 
with time step at 50y.
\item Set the hydraulic conductivity to a low value 
(Spatial attributes/modflow/lpf.8 background value) 
to 0.1 m/y (however there is no head difference so there 
will be no flow so this value plays no role).
\item Run the flow model, the head shall be equal to 1 everywhere 
and for all time steps.
\end{itemize}

\subsection{Transport Characteristics}

\begin{itemize}
\item Set the advection type to TVD 
(in Parameters/Transport/Parameters/Adv.1)
\item Set the diffusion coefficient (will be for all species 
at the beginning), in the dispersion panel 
(Parameters/Transport/Parameters/Dsp.5). 
The diffusion coefficient is given in m$^2$.s$^{-1}$, 
you need to transform this value to be consistent with the models units.
\item Set the porosity value to the regions where the rocks are. 
You can set a 0.05 value for the background 
(Spatial attributes/Mt3dms/Btn.11 
then Background and OK) and create a zone for the bentonite value
\end{itemize}

\subsection{Reactions}
\begin{itemize}
\item Import the database.
\item Open the chemistry dialog (Parameters/Chemistry/Chemistry) 
and go to the Solution panel. In that one create two solutions: 
the background, which will be the granite composition (cf table below) 
and the solution that will be the one in the bentonite. 
Note that the Granite contains quite saline water.
\item Then assign the solution 1 to a zone : 
in spatial attributes/Pht3d/Ph.1 create a zone at the 
place of bentonite (from 5.4 to 6 m) and assign the value 1.0.0.0, 
which means solution 1, mineral, exchange and surface 
assemblage as background (meaning nothing as they were not created). 
The solution 0 (granite) is already set as there is a 0 in the background.
\item Run Pht3D (Parameters/Chemistry/Run), you can have a look at the spatial 
distribution by choosing in Results, a time step and then, 
in Chemistry, a species.
\item In order to see the results go in Parameters/Observation/obs, 
and draw a line along the model. Then in Results/Observation choose profile 
then the name of the line you created.

\item In order to see the importance of the grid, go in the grid, and set the 
cell size (dx) to 25 cm at x=0, 1 cm at x = 5.4 m and 5 cm at 6 m. 
What can you see?
Then add a concentration of 1 mg/L of U(+6), 
which is present when the medium is oxidized. 
Due to its size UO2 diffuse less rapidly than the other ions, its diffusion 
in clays being estimated as 2.6 10-12 m.s-1 
(or five times less than the other ions). 
In order to set a different diffusion coefficient 
for one species, go in Parameters/Transport/Parameters/Dsp.1 
and replace the �#� sign by, on one line, the name of 
the species (�U(+6)   �) followed by its diffusion coefficient 
(using the correct units). 
Look at the results by comparing this ion to the diffusion of the other ions.
\end{itemize}

\begin{table}[hb]
\caption{Aqueous concentrations of the ambient groundwater (Background solution)
and of the tailings fluid (Solution 1).}
\begin{center}
\begin{tabular}{l c c}
\hline
 Aqueous component       & Granite Region           & Bentonite Region        \\ \hline
 Al                      & 2.97 $\times$ $10^{-9}$  & 5.84 $\times$ $10^{-9}$ \\
 C(4)                    & 1.01 $\times$ $10^{-3}$  & 2.27 $\times$ $10^{-3}$ \\
 Ca                      & 5.54 $\times$ $10^{-2}$  & 1.20 $\times$ $10^{-3}$ \\ 
 Cl                      & 1.37 $\times$ $10^{-1}$  & 3.11 $\times$ $10^{-4}$ \\
 Fe(2)                   & 0                        & 0                       \\
 Fe(3)                   & 0                        & 0                       \\
 K                       & 1.68 $\times$ $10^{-4}$  & 0                       \\
 Mg                      & 2.50 $\times$ $10^{-3}$  & 0                       \\
 Na                      & 8.53 $\times$ $10^{-2}$  & 1.00 $\times$ $10^{-4}$ \\
 S(6)                    & 3.19 $\times$ $10^{-2}$  & 6.00 $\times$ $10^{-5}$ \\
 S(-2)                   & 0                        & 0                       \\ 
 Si                      & 1.98 $\times$ $10^{-3}$  & 1.95 $\times$ $10^{-3}$ \\
 Sr                      & 2.93 $\times$ $10^{-4}$  & 0                       \\
 U(6)                    & 0                        & 1.00 $\times$ $10^{-4}$ \\
 U(4)                    & 0                        & 0                       \\
 pH                      & 8.626                    & 7.65                    \\
 pe                      & 8.00 $\times$ $10^{-7}$  & 8.5                     \\ \hline
\end{tabular}
\label{ex_diff_wq}
\end{center}
\end{table}

	Granite	Benthonite
Al		5.84E-09
C(4)	1.01E-03	2.27E-03
Ca	5.54E-02	1.20E-03
Cl	1.37E-01	3.11E-04
Fe(2)	0.00E+00	0.00E+00
Fe(3)	0.00E+00	0.00E+00
K	1.68E-04	0
Mg	2.50E-03	1.20E-04
Na	8.53E-02	1.00E-04
S(6)	3.19E-02	6.00E-05
Si	1.98E-03	1.95E-03
Sr	2.93E-04	0
U(4)	0.00E+00	0.00E+00
U(6)	0.00E+00	1.00E-04
pH	8.626	7.65
pe	11.6	8.5

\subsection{elements diffusion at the boundary in presence of minerals}
 
The first step has been done without minerals, 
which is not correct as the chemical composition is influenced 
by the presence of minerals and some reactions between the solution 
and the minerals. We will thus use Figure 1 to set the content in minerals. 
For this exercise we will keep the solution at equilibrium with the mineral. 
To specify the mineral at equilibrium with the solution:

\begin{itemize}
\item Go into Paramters/Chemistry/Chemistry and in the panel 
Phases tick all the phase that are given in Figure 1. Fe(OH)3(a) 
is present in small amount (0.1 mol/L) Kaolinite and Al(OH)3(a) are 
not present but allowed to precipitate.
\item In the back_SI and Ass1_SI leave 0 everywhere (for equilibrium) 
\item Set the number of moles of mineral for the concerned ones. 
As it is at equilibrium, and the amount of mineral will 
not play a significant role, just start by setting each 
mineral phase amount to 1 mol/L solution. 
The background will be the granite and the assemblage 1, 
will be benthonite.
\item In the spatial attributes/Pht3d/Ph.3 leave the background at 0, 
this will be both the solution 0 and assemblage 0. 
For the bentonite zone already created, go in Modify zone panel, 
select the zone and click on its value to open the dialog and change 
the value from 1.0.0.0 to 1.1.0.0, meaning that the initial solution 
is at equilibrium with mineral assemblage 1 and that these mineral 
will remain in the zone during the simulation.
\item In the solution add Fe(2), Fe(3) and U(4) at a zero concentration, 
\end{itemize}

What did the addition of minerals changed to the concentration profiles, and pH?

The complete case would include the potential precipitation of other minerals, 
but also the addition of ion exchange on the montmorillonite 
and its role on major ions, and the surface complexation of uranium. 


